\section{Actual images from an attained rank-one nilpotent}\label{sec:even-prelim}
The remainder of the article deals with $s_{2r}$. Throughout this part,
its trace-zero property follows directly from cyclic alternation:
for any word of length $2r$, moving the first letter to the last
preserves the trace but changes the permutation sign by
$(-1)^{2r-1}=-1$. Consequently
\begin{equation}\label{eq:even-trace}
 \tr s_{2r}(A_1,\ldots,A_{2r})=0,
 \qquad \im s_{2r}(\Mat_n)\subseteq\Sl_n.
\end{equation}
The identity $s_{2n}=0$ on $\Mat_n$ implies the entire zero branch
for $n\le r$. Moreover $s_{2r}(A_1,\ldots,A_{2r-1},I_n)=0$.
For nonzero images the following local-to-global mechanism is essential.
It is the smooth-map formulation of the nilpotent argument in
\cite[Lemma~2]{Tang2026FiveSix}; the elementary degeneration is
included to make the actual-value conclusion explicit.

\begin{lemma}[Attained-nilpotent actual-image criterion]
\label{lem:nilpotent-open}
Let $F:\mathbb A^m_K\to\Sl_n(K)$ be a polynomial map whose image
contains zero and is invariant under similarities and multiplication
by nonzero scalars. Suppose that for some $x$ the value $R=F(x)$
is nonzero, rank one and square zero, and
\(dF_x:K^m\to\Sl_n(K)\) is surjective. Then
$\im F=\Sl_n(K)$.
\end{lemma}
\begin{proof}
Both affine spaces are smooth. Surjectivity of the differential
makes $F$ smooth at $x$, and smoothness is open; the restriction to
a suitable neighborhood maps onto a Zariski-open subset containing
$R$. In particular, the image contains an open neighborhood $U$ of
$R$ in $\Sl_n$.

Let $T\in\Sl_n$ be nonzero. It is nonscalar in characteristic zero.
A nonscalar linear transformation has a vector $v$ for which
$v,Tv$ are linearly independent: otherwise every line is invariant,
forcing the transformation to be scalar. Thus after a similarity
$T_{01}\ne0$. For $t\ne0$, write
\begin{equation}\label{eq:nilpotent-degeneration}
 D(t)=\diag(t^{-1},t,1,\ldots,1),\qquad
 T(t)=t^2D(t)TD(t)^{-1}.
\end{equation}
All entries of $T(t)$ have nonnegative powers of $t$; its value at
$t=0$ is $T_{01}E_{01}$. Every nonzero rank-one square-zero matrix
is similar, up to a nonzero scalar, to $R$. By similarity and
scaling invariance there is therefore an open neighborhood of
$T_{01}E_{01}$ contained in the image. Its preimage under the
affine-line morphism $t\mapsto T(t)$ contains zero and hence some
nonzero $t$. Thus $T(t)$ is an actual value. Undoing the
similarity and scalar $t^2$ shows that $T$ is an actual value.
The zero target is already in the image.
\end{proof}

A standard polynomial is equivariant for simultaneous conjugation;
scaling one input scales its output, and setting an input to zero
gives zero. Thus every standard evaluation map satisfies the
invariance assumptions of Lemma~\ref{lem:nilpotent-open}.

\section{A uniform minimal-size seed}\label{sec:even-seed}
The first matrix order at which $s_{2r}$ need not vanish is
$n=r+1$. The following construction treats every such order with
one matrix-unit formula, including the full differential required
by Lemma~\ref{lem:nilpotent-open}.

\begin{theorem}[Uniform even-standard seed]\label{thm:even-seed}
For every $n\ge2$,
\begin{equation}\label{eq:even-seed}
 \im s_{2n-2}(\Mat_n(K))=\Sl_n(K).
\end{equation}
More precisely, the evaluation constructed below has a nonzero
rank-one square-zero output and surjective full differential.
\end{theorem}
\begin{proof}
Write $p=n-2$ and index matrix units by $0,\ldots,n-1$.
For $n\ge3$, use the ordered $2p+2$ inputs
\begin{equation}\label{eq:even-tuple}
\begin{aligned}
 \mathcal A=(&L=E_{00},\ X=E_{01},\\[-2pt]
 &S_2=E_{02},T_2=E_{20},\ldots,\\[-2pt]
 &S_{n-1}=E_{0,n-1},T_{n-1}=E_{n-1,0}).
\end{aligned}
\end{equation}
The signed matrix-unit trail calculation gives
\begin{equation}\label{eq:even-seed-evals}
\begin{aligned}
 s_{2n-2}(\mathcal A)&=(p+1)!E_{01},\\
 \partial_{X\gets E_{10}}s_{2n-2}(\mathcal A)&=-(p+1)!E_{10},\\
 \partial_{X\gets E_{12}}s_{2n-2}(\mathcal A)&=-p!E_{12},\\
 \partial_{S_2\gets E_{10}}s_{2n-2}(\mathcal A)&=-(p+1)!E_{20},\\
 \partial_{S_2\gets E_{12}}s_{2n-2}(\mathcal A)&=p!(E_{11}-E_{22}).
\end{aligned}
\end{equation}
The notation $\partial_{A_j\gets V}$ denotes the directional
derivative in the $j$th argument; by multilinearity it is exactly
the value with that argument replaced by $V$.

Here is a signed-trail derivation valid for all $n$. A nonzero
matrix-unit word is precisely an oriented Euler trail on the
labelled directed graph. In the unmodified tuple, each trail starts
at $0$, traverses the loop $L$ and the $p$ excursions
$S_jT_j$ in any order, and ends with $X$. The excursions have
even length and commute in permutation parity; moving $X$ to the
end crosses $2p$ labels. These $(p+1)!$ trails all have positive
sign. With $X=E_{10}$ the trail instead begins with $X$, making
its sign negative. With $X=E_{12}$ it starts with $X,T_2$,
traverses the remaining root blocks and ends with $S_2$;
there are $p!$ such trails and the word has odd inversion parity.
With $S_2=E_{10}$ the trail begins with $T_2$ and traverses the
root loop, the new excursion $XS_2$, and the remaining
$p-1$ excursions; the resulting $(p+1)!$ trails are negative.

Finally, for $S_2=E_{12}$, there is a directed triangle
$X,S_2,T_2$ together with a loop and the other two-edge
excursions. Trails starting at $0$ or another leaf cancel by
exchanging the odd-length loop and triangle. Starting at $1$
leaves $p!$ positively signed trails, whereas starting at $2$
leaves $p!$ negatively signed trails. This proves the last
formula in \eqref{eq:even-seed-evals}, including its diagonal
signs.

Let $\mathcal D$ be the image of the \emph{full} differential of
$s_{2n-2}$ at $\mathcal A$. It lies in $\Sl_n$ by
\eqref{eq:even-trace}. Equivariance supplies every commutator
$[Z,E_{01}]\in\mathcal D$ (the scalar $(p+1)!$ is nonzero).
Consider a change of basis $g=\diag(1,1,G)$,
$G\in\mathrm{GL}_{p}(K)$, acting on the leaf coordinates.
It fixes $L$ and $X$ and sends the $p$ outgoing and $p$ incoming
leaf arguments to mutually contragredient linear combinations
of their original arguments. The combined transformation on
these $2p$ frozen slots has determinant one. The alternating
multilinear map, and hence its differential image at the tuple,
is therefore invariant under conjugation by $g$.

The orbits of $E_{12}$ and $E_{20}$ in
\eqref{eq:even-seed-evals} supply all $E_{1j}$ and $E_{i0}$
for $i,j\ge2$. The orbit of $E_{11}-E_{22}$ spans
\begin{equation}\label{eq:even-seed-block}
 \bigl\{(\tr B)E_{11}-B:
 B\text{ is supported on the leaf block}\bigr\},
\end{equation}
because the conjugates of a rank-one diagonal idempotent span
its full matrix block: rank-one idempotents include $E_{ii}$
and $E_{ii}+E_{ji}$ for distinct leaf indices.
Together with $E_{10}$ and $[Z,E_{01}]$, these directions give
every off-diagonal matrix unit and every diagonal difference.
Thus $\mathcal D=\Sl_n$, as required.

When $n=2$, take $(E_{00},E_{01})$.
Its commutator value is $E_{01}$; replacing $E_{01}$ by $E_{10}$
gives $-E_{10}$ and the conjugation direction gives
$E_{00}-E_{11}$. The full differential is $\Sl_2$.
In each order the output is a nonzero rank-one square-zero
matrix, so Lemma~\ref{lem:nilpotent-open} proves
\eqref{eq:even-seed}.
\end{proof}

\begin{remark}[The size-four seed in Tang's construction]\label{rem:tang-seed}
At $n=4$, the tuple in \eqref{eq:even-tuple} is
$(E_{00},E_{01},E_{02},E_{20},E_{03},E_{30})$.
Qiyue Tang's Appendix~A uses
$(4E_{11},E_{12},E_{13},E_{14},E_{31},E_{41})$ in one-based
coordinates \cite[Proposition~11]{Tang2026FiveSix}.
The two coincide after relabelling the indices, interchanging
the fourth and fifth slots, and multiplying the first slot by
$4$; alternating multilinearity accounts for the sign and
scale. Thus the size-four instance agrees, up to these changes,
with Tang's matrix-unit construction.
\end{remark}

\section{Admissible cores and an additive transfer}\label{sec:even-cores}
The seed can be transferred to larger orders by a normal-map
criterion that keeps track of the \emph{full differential}, rather
than merely of a frozen generic kernel. Fix $r\ge1$ and write
\begin{equation}\label{eq:even-rk}
 k=2r-1.
\end{equation}
For $B=(B_1,\ldots,B_k)\in\Mat_N(K)^k$ and
$u\in K^N$, define
\begin{equation}\label{eq:core-definitions}
\begin{aligned}
 S&=s_k(B_1,\ldots,B_k),\\
 L(P)&=s_{k+1}(B_1,\ldots,B_k,P),\\
 C_i(P)&=s_k(B_1,\ldots,B_{i-1},P,B_{i+1},\ldots,B_k),\\
 \Gam_{B,u}(P)&=\bigl(L(P),C_1(P)u,\ldots,C_k(P)u\bigr).
\end{aligned}
\end{equation}
Thus $\Gam_{B,u}:\Mat_N(K)\to\Mat_N(K)\oplus(K^N)^k$
is a linear map in $P$; the frozen data $(B,u)$ vary
polynomially. Call the size $N$ \emph{admissible for $r$} if
there exists one such tuple with
\begin{equation}\label{eq:core-admissible}
 \det S\ne0,\qquad \ker\Gam_{B,u}=\{0\}.
\end{equation}
Both conditions are polynomial nonvanishing conditions in
unrestricted parameters. A witness therefore defines a nonempty
Zariski-open subset of the core parameter space.

\begin{lemma}[Normal-map criterion]\label{lem:core-normal}
If $N$ is admissible for $r$, then
\begin{equation}\label{eq:core-conclusion}
 \im s_{2r}(\Mat_{N+1}(K))=\Sl_{N+1}(K).
\end{equation}
Indeed the defining witness yields an attained nonzero rank-one
square-zero value at which the full differential is surjective.
At the particular border tuple used below, assuming $S$ invertible,
this differential condition is equivalent to injectivity of
$\Gam_{B,u}$.
\end{lemma}
\begin{proof}
Embed the fixed matrices and the final input as
\begin{equation}\label{eq:core-border}
 A_i=\begin{pmatrix}B_i&0\\0&0\end{pmatrix},\qquad
 Y=\begin{pmatrix}0&u\\0&0\end{pmatrix}.
\end{equation}
Only words ending in $Y$ can survive, because $YA_i=0$.
Since $k$ is odd, the value of the even alternator is
\begin{equation}\label{eq:core-border-value}
 s_{k+1}(A_1,\ldots,A_k,Y)
 =\begin{pmatrix}0&Su\\0&0\end{pmatrix}.
\end{equation}
If $u=0$, then $\Gam=L$ and its kernel contains all
$B_i$ (and also the identity), contradicting injectivity.
Hence $Su\ne0$. The value is nonzero, square zero and rank one.

To compute the differential, pair its values with a general
matrix using the nondegenerate trace form. Variations of $Y$
in the upper and lower borders give the full upper and lower
border directions by left and right multiplication by $S$ and
$-S$, respectively. Thus every normal is block diagonal.
The differential values have trace zero, so scalar matrices
are normals by \eqref{eq:even-trace}. Subtracting a scalar we
may write a normal uniquely as $\diag(P,0)$.

Variations of the old $N\times N$ block of $Y$ give
$\tr(P L(Q))$ for arbitrary $Q$. The map $L$ is
skew-adjoint for the trace pairing. To see this, collect the
cyclic classes in
$\tr s_{k+2}(B_1,\ldots,B_k,Q,P)$. The length
$k+2$ is odd, so a cyclic shift preserves sign; moving $Q$
past the $k$ frozen letters contributes a minus sign.
Consequently the trace equals $-(k+2)\tr(Q L(P))$.
Interchanging $P,Q$ in the alternating trace yields
\begin{equation}\label{eq:core-skew}
 \tr(P L(Q))=-\tr(Q L(P)).
\end{equation}
The old-block normal equations are therefore $L(P)=0$.

Now vary the $i$th $A_i$ in the lower border by a row
$c\in(K^N)^*$. For a nonzero word the unique upper-border
factor $Y$ must be immediately followed by this lower-border
variation. Remove $Y$ and rotate the remaining $k$ letters
cyclically. If $a$ of the original old letters followed the
adjacent pair, the removal has sign $(-1)^{a+1}$ and moving
the varied letter to its original position has sign $(-1)^a$.
The product is $-1$, and the normal pairing is exactly
\begin{equation}\label{eq:core-low-normal}
 -c\,C_i(P)u.
\end{equation}
Because all rows $c$ are available, these normal equations say
$C_i(P)u=0$ for each $i$. Therefore every normal modulo
scalars satisfies $\Gam_{B,u}(P)=0$. Conversely, in the pure
border tuple \eqref{eq:core-border}, these are all the normal
equations: other source-block variations yield only upper-border
outputs or zero. This also proves the asserted equivalence.
If $\Gam$ is injective, all normals are scalar, so the
full differential has image $\Sl_{N+1}$.
Lemma~\ref{lem:nilpotent-open} now proves
\eqref{eq:core-conclusion} as an \emph{attained} image equality.
\end{proof}

A complementary rectangular nonvanishing calculation records the
invertibility available for independently varying block-diagonal inputs.
\begin{lemma}[Mixed-rectangle nonvanishing]\label{lem:mixed-rect}
Let $a,b\ge1$ and $m<2(a+b)$. For an $m$-variable standard
polynomial, the free rectangular map
$\Mat_{a\times b}\to\Mat_{a\times b}$ obtained by freezing
$m-1$ block-diagonal inputs of sizes $a$ and $b$ is generically
invertible as the frozen blocks vary independently.
\end{lemma}
\begin{proof}
Choose integers $e,f\ge0$ with $e+f=m-1$,
$e<2a$ and $f<2b$; they exist since $m<2(a+b)$.
Specialize the first $e$ frozen labels to the $a$ block and
the remaining $f$ to the $b$ block. The only surviving words
put the left labels before the free rectangle and the right
labels after it, so the map is
\begin{equation}\label{eq:mixed-rect}
 X\longmapsto(-1)^f s_e(A)X s_f(B),\qquad s_0=I.
\end{equation}
Both factors admit invertible evaluations. Indeed, for degree
below the Amitsur--Levitzki threshold a standard polynomial
is a nonidentity. Expand a nonzero evaluation in the
clock-and-shift homogeneous basis of the corresponding matrix
algebra; by multilinearity one homogeneous tuple yields a
nonzero scalar multiple of an invertible basis element.
Thus the rectangular map is invertible at one specialization,
and its determinant is a nonzero polynomial.
\end{proof}

\begin{theorem}[Additivity of admissible core sizes]\label{thm:core-add}
For a fixed $r\ge1$, if $a,b\ge r$ are admissible core sizes,
then $a+b$ is admissible for $r$.
\end{theorem}
\begin{proof}
Choose witnesses $(B^{(a)},u_a)$ and $(B^{(b)},u_b)$
for the two sizes. For one parameter $t$ set
\begin{equation}\label{eq:core-add-tuple}
 B_i(t)=\diag(B_i^{(a)},tB_i^{(b)}),\qquad
 u=\binom{u_a}{u_b}.
\end{equation}
For every $t\ne0$, the diagonal endpoint blocks are
$S_a$ and $t^kS_b$ and hence are invertible. The old-block
$L$ and $C_i$ operators in the second block are respectively
scaled by $t^k$ and $t^{k-1}$, so their individual
admissibility kernels remain zero.

On the two off-diagonal rectangular spaces, the map $L_t$
in \eqref{eq:core-definitions} is polynomial in $t$. At $t=0$
the top-right rectangular map is $X\mapsto S_aX$,
because only words with the free variable last survive.
The bottom-left rectangular map is $X\mapsto -XS_a$,
because only words with the free variable first survive and
$k$ is odd. Both are invertible. Therefore their determinants
are nonzero polynomials in the \emph{same} scalar parameter.
As $K$ is infinite, select one common $t\in K^\times$
avoiding their zero sets.

At this fixed $t$, if $\Gam_{B(t),u}(P)=0$, the two mixed
$L_t$ equations force the off-diagonal blocks of $P$ to be
zero. The remaining $C_i(P)$ are block diagonal, and the
normal-map equations reduce to those for the two original
cores, with nonzero scalar factors on the second block.
Both diagonal blocks of $P$ vanish by the original
injectivity hypotheses. Thus $\Gam_{B(t),u}$ is injective,
and $\det s_k(B(t))\ne0$. This constructs a single
admissible witness of size $a+b$.
\end{proof}

\begin{corollary}[An infinite congruence family]\label{cor:congruence}
For every $r,q\ge1$,
\begin{equation}\label{eq:congruence}
 \im s_{2r}(\Mat_{qr+1}(K))=\Sl_{qr+1}(K).
\end{equation}
\end{corollary}
\begin{proof}
In Theorem~\ref{thm:even-seed}, set $n=r+1$, place the
exit variable last, and remove its sink coordinate. The old
$r\times r$ core has
\begin{equation}\label{eq:core-star}
 B=(E_{00},E_{01},E_{10},\ldots,E_{0,r-1},E_{r-1,0}),
 \qquad u=e_0,
\end{equation}
with the obvious singleton interpretation when $r=1$.
Counting closed trails through the root loop and the $r-1$
two-edge excursions gives
\begin{equation}\label{eq:core-star-endpoint}
 S=s_{2r-1}(B)
 =\diag(r!,(r-1)!,\ldots,(r-1)!),
\end{equation}
which is invertible in characteristic zero.
The full differential in Theorem~\ref{thm:even-seed} is onto
$\Sl_{r+1}$; the converse statement in
Lemma~\ref{lem:core-normal} shows that this old core is
admissible. For $r=1$ it follows also directly from
$S=1$ and $C_1(P)u=P$. Add $q$ copies of the admissible
size $r$ with Theorem~\ref{thm:core-add}, then apply
Lemma~\ref{lem:core-normal} in size $qr+1$.
\end{proof}

\section{Exact core certificates in degrees eight and ten}\label{sec:even-low}
For each $r$, the uniform construction supplies an admissible core
of size $r$. By additive transfer, to obtain all sizes $N\ge r$,
it suffices to establish admissibility for the finite range
\begin{equation}\label{eq:core-finite-reduction}
 r,r+1,\ldots,2r-1.
\end{equation}
Indeed every $N\ge r$ has the form $N=qr+s$,
$q\ge1$ and $0\le s<r$. If $s=0$, use $q$ copies of
$r$; otherwise use $q-1$ copies of $r$ and one copy of $r+s$.
Thus a finite collection of admissible cores extends to every
larger core order by Theorem~\ref{thm:core-add}.

For $r=4$ and $r=5$, the seven residual cases have exact
integer witnesses. For each case in the table, the corresponding
certificate contains an explicit ordered tuple $B$ of integer matrices,
an integer column $u$,
$\det S\ne0$, and the row indices and exact nonzero integer
determinant of an $N^2\times N^2$ submatrix of
$\Gam_{B,u}$. The last two columns give reductions of these full integer
minors modulo $101$ and $103$; their nonzero residues verify
nonvanishing in characteristic zero.
\begin{center}
\small
\begin{tabular}{@{}rrrcrr@{}}
\toprule
$r$ & $N$ & $2r$ & $\rk\Gam$ & minor $\bmod 101$ & minor $\bmod 103$\\
\midrule
4&5&8&25&48&57\\
4&6&8&36&36&25\\
4&7&8&49&75&58\\
5&6&10&36&52&34\\
5&7&10&49&64&45\\
5&8&10&64&45&60\\
5&9&10&81&85&7\\
\bottomrule
\end{tabular}
\end{center}
The certificate names are listed in
Appendix~\ref{app:even-verification}; they encode each pair $(r,N)$
from the table. For example,
\path{certificates/border_r4_N5.json} records $(r,N)=(4,5)$. The normal map has $N^2$
columns, so a nonzero $N^2$-minor proves its injectivity.
The invertible endpoint is separately recorded for each case.

\begin{theorem}[Degrees eight and ten without generic-nullity input]
\label{thm:even-eight-ten}
For every algebraically closed characteristic-zero field $K$,
\begin{equation}\label{eq:eight-ten}
\begin{aligned}
 \im s_8(\Mat_n(K))&=
 \begin{cases}\{0\},&n\le4,\\\Sl_n(K),&n\ge5,\end{cases}\\
 \im s_{10}(\Mat_n(K))&=
 \begin{cases}\{0\},&n\le5,\\\Sl_n(K),&n\ge6.\end{cases}
\end{aligned}
\end{equation}
\end{theorem}
\begin{proof}
\enlargethispage{2\baselineskip}
For $r=4$, the uniform seed gives admissible size $N=4$,
and the three finite residual sizes in
\eqref{eq:core-finite-reduction} are $N=5,6,7$.
For $r=5$, the uniform seed gives $N=5$, and the finite
residual sizes are $N=6,7,8,9$. Thus the seven exact
certificate witnesses together with Theorem~\ref{thm:core-add}
make every core size $N\ge r$ admissible, for $r=4,5$.
Lemma~\ref{lem:core-normal} yields the actual traceless image
for every order $n=N+1\ge r+1$. The other orders vanish by
Amitsur--Levitzki.
\end{proof}

Appendix~\ref{app:even-verification} describes how to check the
seven saved certificates. The extension to all matrix orders
uses Theorem~\ref{thm:core-add}.

\section{Generic nullity and the all-even classification}\label{sec:even-general}
We now deduce the remaining even degrees from two precise
external generic-kernel theorems. The argument extends the
approach to $s_6$ in Qiyue Tang's Appendix~A
\cite[Proposition~11]{Tang2026FiveSix}; in particular,
its use of a rank-one nilpotent, a generic central division
algebra and a projective eigenline incidence is directly indebted
to that source.

Write
\begin{equation}\label{eq:generic-op}
 \mathcal L_{A_1,\ldots,A_j}(P)
 =s_{j+1}(A_1,\ldots,A_j,P).
\end{equation}
The two inputs needed below are as follows. Brassil and
Reichstein \cite[Theorem~1.2]{BrassilReichstein2022} prove that
for $j$ even, $2\le j\le2N-2$, and generic $j$-tuples in
$\Mat_N$, the kernel of \eqref{eq:generic-op} has dimension
$j$ in characteristic zero. Their theorem is stated more
generally over an infinite field with specified characteristic
exclusions; only the characteristic-zero specialization is
used here. Lin, Chenhao Tang and Zhao
\cite[Theorem~1.2 and Remark~6.1]{LinTangZhao2026}
prove the complementary odd-frozen statement: for $N\ge3$,
odd $j$ with $3\le j\le2N-3$ and any characteristic-zero
field, the generic nullity is
\begin{equation}\label{eq:GN}
 \dim\ker\mathcal L_{A_1,\ldots,A_j}
 =\begin{cases}j+1,&N\text{ even},\\
 j+2,&N\text{ odd}.\end{cases}
\end{equation}
These generic-kernel dimensions provide the input for the
following attained-image argument.

\begin{theorem}[Even-standard classification from generic nullity]
\label{thm:even-complete}
For every algebraically closed field $K$ of characteristic zero
and all $r,n\ge1$,
\begin{equation}\label{eq:even-complete}
 \im s_{2r}(\Mat_n(K))=
 \begin{cases}\{0\},&n\le r,\\\Sl_n(K),&n\ge r+1.\end{cases}
\end{equation}
For $r\ge2$ and $n\ge r+2$, the proof uses the above external
Brassil--Reichstein and Lin--Tang--Zhao generic-nullity theorems.
\end{theorem}
\begin{proof}
If $n\le r$, then $2r\ge2n$, and the standard identity
forces the image to be zero. The case $r=1$ is the classical
commutator image theorem \cite{AlbertMuckenhoupt1957}.
For $r\ge2$ and $n=r+1$, the positive claim is
Theorem~\ref{thm:even-seed}.

Now let $r\ge2$ and $n\ge r+2$. Set
\begin{equation}\label{eq:generic-parameters}
 N=n-1,\qquad k=2r-1;
 \quad N\ge r+1\ge3,\quad 3\le k\le2N-3.
\end{equation}
For each $C_i$ of \eqref{eq:core-definitions}, the number
$k-1=2r-2$ of frozen inputs is even and satisfies
$2\le k-1\le2N-2$. Therefore both external nullity theorems
apply in precisely their stated ranges.

Take $k$ independent generic $N\times N$ matrices
$B_1,\ldots,B_k$ over $K$, and let $D$ denote the central
division algebra of generic matrices generated by them, with
center $F$ and degree $N$. Write $E$ for the field of their
independent matrix-entry indeterminates. The natural generic
matrix representation embeds $D$ in $\Mat_N(E)$ and splits it
over $E$; the center $F$ embeds in $E$ as scalar matrices.
The matrices $B_i$ are linearly independent over $F$,
because their vectorized columns are independent already over
$E$ (some $k\times k$ minor is a nonzero polynomial, as
$k\le N^2$). Set
\begin{equation}\label{eq:generic-S}
 S=s_k(B_1,\ldots,B_k)\in D.
\end{equation}
Since $k<2N$, the standard polynomial is not an identity
of $\Mat_N$; genericity makes $S\ne0$ and thus
$S\in D^\times$. In the generic matrix model, all of the
maps $L,C_i$ from \eqref{eq:core-definitions} are
$F$-linear operators on the $N^2$-dimensional space $D$.
Their ranks persist under splitting-field extension.

The Brassil--Reichstein theorem gives, simultaneously for
every $i$,
\begin{equation}\label{eq:generic-C-kernel}
 \ker C_i=\Span_F\{B_j:j\ne i\}.
\end{equation}
Indeed alternation gives the displayed $(k-1)$-dimensional
subspace, and the generic theorem gives exactly this nullity.
The generic open conditions for the finitely many $C_i$
intersect in the irreducible generic input space, so one set
of generic entries realizes all the ranks at once.
The odd-frozen theorem \eqref{eq:GN} gives
\begin{equation}\label{eq:generic-W}
 W=\ker L/\Span_F(B_1,\ldots,B_k),\qquad
 \dim_F W=\begin{cases}1,&N\text{ even},\\2,&N\text{ odd}.
 \end{cases}
\end{equation}
(The upper bound $\dim W\le2$ would suffice below.)

Alternation also gives
\begin{equation}\label{eq:C-on-B}
 C_i(B_j)=\delta_{ij}S.
\end{equation}
The $F$-linear maps
\begin{equation}\label{eq:generic-T}
 T_i:W\longrightarrow D/F1,
 \qquad[P]\longmapsto S^{-1}C_i(P)+F1
\end{equation}
are consequently well defined. They are \emph{jointly
injective}. In fact, if $C_i(P)=b_iS$ for every $i$, then
\eqref{eq:C-on-B} gives
$P-\sum_i b_iB_i\in\bigcap_i\ker C_i=\{0\}$;
the last equality follows from \eqref{eq:generic-C-kernel}
and the linear independence of all the $B_i$.
This joint injectivity persists after scalar extension.

Choose an $F$-linear section $W\to\ker L$ and let
$\bar F$ be an algebraic closure of $F$. Fix a splitting
identification $D\otimes_F\bar F\simeq\Mat_N(\bar F)$.
For a class $[P]\in W_{\bar F}$ the endomorphisms
$S^{-1}C_i(P)$ are defined by this section; changing the
section adds scalars to them, so their eigenlines are
independent of the section. We claim that there is a nonzero
column $u\in\bar F^N$ which is not a common eigenvector of
$S^{-1}C_i(P)$ for \emph{any} nonzero class $[P]$.

If $\dim W=0$ the claim is immediate. If $\dim W=1$, joint
injectivity means at least one of these matrices is nonscalar;
its eigenvectors form a finite union of proper eigenspaces,
which does not exhaust $\bar F^N$. The substantive case is
$\dim W=2$. In that case form the projective closed incidence
\begin{equation}\label{eq:generic-incidence}
 \mathcal I=\bigl\{([P],[u])\in
 \mathbb P(W_{\bar F})\times\mathbb P^{N-1}_{\bar F}:
 u\wedge S^{-1}C_i(P)u=0\quad(1\le i\le k)\bigr\}.
\end{equation}
The incidence records \emph{simultaneously} all kernel classes
and their common eigenlines.

At the generic point of the projective line
$\mathbb P(W_{\bar F})$, write $P=P_0+tP_1$ using an
$F$-basis of $W$. Joint injectivity of \eqref{eq:generic-T}
shows that some $S^{-1}C_i(P)$ is noncentral in
$D\otimes_FF(t)$. This algebra is a division algebra:
it embeds in the Laurent-series division ring $D((t))$,
with central indeterminate $t$, and as a finite-dimensional
domain over $F(t)$ every nonzero element is invertible.
A noncentral element in that division algebra has an
irreducible separable minimal polynomial of degree $a\ge2$
over $F(t)$; its degree divides $N$. Upon splitting, its
reduced characteristic polynomial is the $N/a$-th power of
that minimal polynomial. Its distinct eigenvalues all have
eigenspace dimension $N/a\le\lfloor N/2\rfloor$.
Consequently the generic fiber of \eqref{eq:generic-incidence}
has dimension at most $\lfloor N/2\rfloor-1$.

At every special geometric point $[P]$, the scalar-extended
joint injectivity still ensures that some
$S^{-1}C_i(P)$ is nonscalar. Its eigenspaces have dimension
at most $N-1$, and hence the corresponding fiber has
dimension at most $N-2$. Thus any component of
$\mathcal I$ dominating the projective line has dimension
at most $\lfloor N/2\rfloor$, and a component over a
point has dimension at most $N-2$. Since $N\ge3$,
\begin{equation}\label{eq:incidence-bound}
 \dim\mathcal I\le N-2<N-1=\dim\mathbb P^{N-1}.
\end{equation}
The projection to $\mathbb P^{N-1}$ is proper, so its image
is a proper closed subset. Choose $[u]$ outside that image.
This proves the claim for the \emph{same} $u$ and every
nonzero kernel class.

For this column, the normal map $\Gam_{B,u}$ is injective
after the chosen splitting. Indeed $\Gam(P)=0$ first implies
$P\in\ker L$. If its class in $W$ were nonzero, all
$C_i(P)u=0$ would make $[u]$ a common eigenline, contrary
to construction. Therefore
$P=\sum_i b_iB_i$; then
$C_i(P)u=b_iSu=0$ for all $i$ forces each $b_i=0$.
We have also retained $\det S\ne0$.

It remains to obtain a witness over the originally specified
field $K$, rather than merely over an abstract splitting of
$D$. Choose an algebraic closure $\Omega$ of the
original generic-entry field $E$ containing an $F$-embedding
of $\bar F$. The generic matrix representation and the
chosen split representation of
$D\otimes_F\Omega\simeq\Mat_N(\Omega)$ differ by an inner
automorphism, by Skolem--Noether. Transport the column $u$
under that conjugacy. Equivariance of $S,L,C_i$ shows that
$\det S\ne0$ and $\rk\Gam=N^2$ hold in the original
generic-entry representation over $\Omega$.

\enlargethispage{1\baselineskip}
Both assertions mean that the product of $\det S$ with a
suitable $N^2\times N^2$ minor of $\Gam$ is a \emph{nonzero
polynomial} in the generic entries of $B$ and the coordinates
of $u$, with coefficients in $K$. Because $K$ is infinite,
there is one specialization of all these coordinates to $K$
at which this polynomial is nonzero. This gives an admissible
core of size $N=n-1$ over $K$ itself, in the sense of
\eqref{eq:core-admissible}. By Lemma~\ref{lem:core-normal},
the full polynomial map $s_{2r}$ attains a nonzero rank-one
nilpotent and therefore has actual image $\Sl_n(K)$.
This proves \eqref{eq:even-complete} in the remaining range.
\end{proof}

